The direct branch keeps every quantified coordinate unrestricted over
$\mathbb Z$ and uses only small integer coefficients. Its principal
equation is
\begin{equation}\label{eq:principal}
 \boxed{\displaystyle
 P_{\DirectVars}(x_0,\ldots,x_{\DirectLast})
   =\sum_{j=0}^{\DirectRowLast}r_j(\mathbf x)^2=0,
 \qquad\mathbf x\in\mathbb Z^{\DirectVars}.}
\end{equation}
Every residual is the explicit integer quadratic printed in
Appendix~\ref{app:rows}. The archive also supplies every coefficient of
the fully collected quartic. A JSON term $[c,[i_1,\ldots,i_k]]$ means
$cx_{i_1}\cdots x_{i_k}$; an empty index list means a constant.
The list contains all quantified coordinates. There are no unnamed
existential witnesses.

This construction combines three separately attributed ingredients:
the factorial-free arithmetic source with the affine radix described
below, the strict local ranges and fixed-base pooling of the incoming
Pooled Powers manuscript \cite{pooled}, and the integral quadratic
projection of the incoming 234-variable manuscript \cite{projection}.
The new deliverable is their compatible, fully emitted composition.
The earlier methods and the already submitted diagonal growth choice
are not presented as discoveries of this report.

\subsection{Two quadratic rows certify a block of powers}

The inherited signed Pell frame provides integers $W\ge258$ and $A$
with $|A|>W^W$. For a block of $r\le15$ requests with the same
nonzero integer base $n$, it supplies canonical signed Pell pairs
$(X_i,Y_i)$ at signed indices $k_i$ and coordinates $Z_i=AY_i$.
Every base, proposed output and index has square less than $W$.
Put
\[
 T=2An-n^2-1,\qquad D_i=m_i-X_i+Z_i,
\]
and use the following quantities only in the proof:
\[
 H=\sum_iY_i(m_i-X_i),\qquad
 S=\sum_i(D_i^2-Y_i^2).
\]
Instead of a separate quotient for every requested power, introduce
just two integer coordinates $K,L$, with the quadratic rows
\begin{align}
 K-2\sum_iY_i(m_i-X_i)-(2A-n)L&=0,\label{eq:pool1}\\
 \sum_i\bigl((m_i-X_i+Z_i)^2-Y_i^2\bigr)-nK+L&=0.
 \label{eq:pool2}
\end{align}
No coordinate for $H$, $S$, or $T$ is needed.

Eliminating $K$ in the proof gives $S-2nH=TL$. The exact identity
\[
 \sum_i(D_i-nY_i)^2=S-2nH-T\sum_iY_i^2
\]
therefore makes the left side divisible by $T$.
Let $b=\lfloor\sqrt{W-1}\rfloor$. Then
$|n|,|m_i|,|k_i|\le b$, and the Pell congruence, multiplied by
$n^b$ to clear every negative index, makes $T$ divide
\[
 E_*:=\sum_i\bigl(m_i n^b-n^{b+k_i}\bigr)^2.
\]
All exponents here are nonnegative proof notation. The estimates
\[
 0\le E_*\le4r b^{4b}<W^{2b+1}\le W^W<|T|
\]
force $E_*=0$. Each square is zero. If $|n|\ge2$, a negative
$k_i$ would give the impossible integer identity $m_i n^{-k_i}=1$.
Thus all $k_i\ge0$ and $m_i=n^{k_i}$.

Conversely, genuine powers and canonical Pell pairs make each
$D_i-nY_i$ divisible by $T$. The identity makes
$L=(S-2nH)/T$ integral, and \eqref{eq:pool1} defines integral $K$.
This proves both directions of the pooled certificate. Its proof uses
the existing frame; no larger frame or new variable-exponent operation
has been inserted into the polynomial.

\subsection{Strict local ranges and the sign order}

The incoming strict-horizon argument chooses a cofinal proof horizon
strictly larger than all finitely many local labels and parameters.
Consequently its local binomial range tests may enforce $0\le u<t$
instead of $0\le u\le t$. In the literal source, the local range
denominator is changed from $q^{t+1}$ to the already required $q^t$.
The four now private coordinates and their rows can be removed with
the stated restoration identities. This is an existential completeness
change justified by choosing a larger horizon; it is not an equality
of arbitrary old and new coordinate tuples.

After this change the source has eleven requests at $q+1$ and three at
$q$. The three other requests remain individual. Applying
\eqref{eq:pool1}--\eqref{eq:pool2} to the two blocks gives a literal
223-coordinate system of 160 quadratic residuals. The finite certificate
authenticates the exact affine offsets of all 14 requests, all retained
$AY$ and Pell norm rows, and every base/index/output in the square frame.

In our arithmetic source $q=6z+3$ holds before any sign inference.
Hence $|q|\ge3$ and $|q+1|\ge2$ for every integer $z$. The block
theorem immediately restores all 14 individual powers and their
nonnegative indices, after the inherited orientation-first theorem
has supplied the canonical pairs. Thus the sign bootstrap in the next
sections remains acyclic. Completeness starts from genuine powers and
uses the converse part of the block theorem to construct the four new
pooled coordinates.

\subsection{Integral projection along a common quadratic direction}

Here is the exact algebra behind the final six projections. It also
applies to the separately retained unchanged 225-variable source.

\begin{lemma}[Quadratic projection with a unit slope]\label{lem:projection}
Let $Q_0,\ldots,Q_{r-1}\in\mathbb Z[x_0,\ldots,x_{n-1}]$ have
degree at most two. Suppose $v\in\mathbb Z^n$ is in the common kernel
of their quadratic Hessians, $v_p=1$, and
$\partial_vQ_j=a_j\in\mathbb Z$ for every $j$.
If $a_s=\pm1$, the common integer zero set is in polynomial bijection
with the common zero set in $n-1$ variables of
\[
 R_j(y)=\left.(Q_j-a_sa_jQ_s)\right|_{y_p=0},\qquad j\ne s.
\]
Every $R_j$ still has degree at most two.
\end{lemma}
\begin{proof}
The common-kernel condition gives $Q_j(y+vt)=Q_j(y)+a_jt$.
For $y_p=0$, set $T(y)=-a_sQ_s(y)$ and $x=y+vT(y)$.
All $Q_j(x)$ vanish whenever all $R_j(y)$ vanish.
Conversely, for a common integer zero $x$ put $y=x-vx_p$.
Then $y_p=0$, and $Q_s(x)=0$ forces $x_p=T(y)$.
The two polynomial maps are inverse on the zero sets. No division,
positivity condition, or unproved divisibility is involved.
\end{proof}

The six literal certificates apply Lemma~\ref{lem:projection} to the
223-coordinate source and give \DirectVars\ coordinates and
\DirectRows\ residuals. Their independently written checker obtains
each reduced row directly by the displayed subtraction identity,
rather than invoking the producer's change-of-coordinates code.
It checks the full direction, every slope, the unit condition, every
new coefficient, and the inverse maps.

Finally elementary integer row operations $Q_i\leftarrow Q_i+kQ_j$
reduce coefficient height and select useful support/height tradeoffs.
The inverse uses $-k$ in reverse order. Such changes preserve
simultaneous vanishing, while changing the sum-of-squares polynomial
itself. The selected polynomial is consequently identified by its own
full coefficient file and cryptographic hash.

Because a sum of integer squares vanishes exactly when every residual
vanishes, these explicit zero-set maps and the arithmetic/refutation
calibration prove the equivalence claimed for \eqref{eq:principal}.
